latex tips

% new column!!
    &

% newline line break!!
    \\

% show "box" for end of maths proof
\hfill $\square$\\[1cm]

% Google book bibtext to APA conversion
    http://ref.lexique.org/

% install full latex system
    % ai texlive texlive-latex-extra texlive-latex3 texlive-fonts-extra 
    ai texlive texlive-latex-extra texlive-fonts-extra 
    ai texlive-science
    ai texlive-xetex # (handles utf8 characters)
    % Get Elsevier "elsarticle.cls"
    ai texlive-publishers
    % 'Matplotlib' needs this
    ai cm-super

% missing 'sty' packages
   ai texlive-latex-extra
   ai texlive-latex3
   ai texlive-fonts-extra
   or for my thesis copy "Missing Packages" folder out into root

% include astro journals (define \apj etc)
    \usepackage{aas_macros}

% graphicspath / images

    % SYMLINK 'images/' to where-ever, then:
        \graphicspath{{images/}}
        \graphicspath{{images//}}

    % BASH enviornment variable
        % (but ensure 'graphicspath' is NOT set)
        export TEXINPUTS=$MY_IMAGES/astro//:

    % 'graphicspath'
        % include a FINAL SLASH at end of path!
        % include a FINAL SLASH at end of path!
        \graphicspath{{/home/ef/study/images/astro/}}
        \graphicspath{{\$HOME/study/images/astro/}}
                       ^
                       |-- '\$' causes syntax color problem!
% brackets!
    \big( \Big( \bigg( \Bigg(
    \big] \Big] \bigg] \Bigg]
    \big\{ \Big\{ \bigg\{ \Bigg\{
    \big \langle \Big \langle \bigg \langle \Bigg \langle
    \big \rangle \Big \rangle \bigg \rangle \Bigg \rangle

% display an ampersand in text \"&"
    \&

% better display fractions
    \dfrac (my "df" trigger)

% compress crush make lines closer
    \usepackage{setspace}
    % \setstretch{1.3}
    \setstretch{1.0}

% math bold font (I never knew!)
    .. = ... + \mathbf{ g_{ij}h^2 }

% congruent
    \equiv

% align columns: better control
    \begin{alignat*}{9}
        \text{top row } \qquad & 00000 \qquad &  01011 \qquad \\
    \end{alignat*}

% align environment: tighten reduce vertical space and also space before and after it
    {
    \setlength{\abovedisplayskip}{0pt}
    \setlength{\belowdisplayskip}{0pt}
    \addtolength{\jot}{-0.5em}
    \begin{align*}
    ...
    \end{align*}
    }%

% misc tip
    "Insert a \% symbol after closing the group with \tiny"
    {\tiny
    }%
    Following paragraph.

% user defined paths!
    \newcommand*{\week}{Week 01}
    \newcommand*{\CommonPath}{../common}
    \input{\CommonPath/preamble.tex}%
    \begin{document}
    \input{title.tex}%

% paragraph spacing
    % method 1:
        \setlength{\parskip}{25pt plus 1pt minus 1pt}
        \setlength{\parindent}{0pt}
        % \setlength{\parskip}{6em}
    % method 2:
        \usepackage[parfill]{parskip}

% figures side by side
    \usepackage{subfig}
    \begin{figure}%
        \centering
        \subfloat[1D function]{{\includegraphics[trim=2cm 0 0 0, scale=0.5]{1dfunction.jpg} }}%
        \subfloat[2D function]{{\includegraphics[scale=0.5]{2dfunction.jpg} }}%
        \caption{Difference plotting 1D and 2D functions}%
        \label{fig:functions}%
    \end{figure}

% table span pages
    \usepackage{longtable}
    \begin{longtable}{l|c|c}
    ...
    \end{longtable}

% supress page numbers
    \pagenumbering{gobble}

% horizontal line subtract (\cline{1-4})
    % method 1:
        \[ \arraycolsep=1.4pt\def\arraystretch{1.2}
        \begin{array}{llll}
            010& 011& 110& 101 \\
            011& 011& 101& 000 \ + \\
            \cline{1-4}
            001& 000& 011& 101
        \end{array}
        \]

    % other method:
        \cline{1-2}
        \\[-1.0em] \cline{1-2} \\[-2.0em]

% f prime
    f^\prime

% break nicely for long urls
    \usepackage{hyperref}
    \def\UrlBreaks{\do\/\do-}

% source code listings
    \usepackage{listings}
    \usepackage{color}

    \definecolor{mygreen}{rgb}{0,0.6,0}
    \definecolor{mygray}{rgb}{0.5,0.5,0.5}
    \definecolor{mymauve}{rgb}{0.58,0,0.82}

    % specify path to source code
    \lstset{inputpath="c:\users\path\to\my\source\code"}

    \lstset{ %
    language=Python,                 % the language of the code
    backgroundcolor=\color{white},   % choose the background color; you must add \usepackage{color} or \usepackage{xcolor}
    basicstyle=\footnotesize,        % the size of the fonts that are used for the code
    breakatwhitespace=false,         % sets if automatic breaks should only happen at whitespace
    breaklines=true,                 % sets automatic line breaking
    captionpos=b,                    % sets the caption-position to bottom
    commentstyle=\color{mygreen},    % comment style
    deletekeywords={...},            % if you want to delete keywords from the given language
    escapeinside={\%*}{*)},          % if you want to add LaTeX within your code
    extendedchars=true,              % lets you use non-ASCII characters; for 8-bits encodings only, does not work with UTF-8
    frame=single,	                   % adds a frame around the code
    keepspaces=true,                 % keeps spaces in text, useful for keeping indentation of code (possibly needs columns=flexible)
    keywordstyle=\color{blue},       % keyword style
    otherkeywords={*,...},           % if you want to add more keywords to the set
    numbers=left,                    % where to put the line-numbers; possible values are (none, left, right)
    numbersep=10pt,                   % how far the line-numbers are from the code
    numberstyle=\tiny\color{mygray}, % the style that is used for the line-numbers
    rulecolor=\color{black},         % if not set, the frame-color may be changed on line-breaks within not-black text (e.g. comments (green here))
    showspaces=false,                % show spaces everywhere adding particular underscores; it overrides 'showstringspaces'
    showstringspaces=false,          % underline spaces within strings only
    showtabs=false,                  % show tabs within strings adding particular underscores
    stepnumber=1,                    % the step between two line-numbers. If it's 1, each line will be numbered
    stringstyle=\color{mymauve},     % string literal style
    tabsize=2,	                   % sets default tabsize to 2 spaces
    title=\lstname                   % show the filename of files included with \lstinputlisting; also try caption instead of title
    }

    % paste in some code
    \begin{lstlisting}
        .... my code ...
    \end{lstlisting}

    % or:
    % specify an actual file
    \lstinputlisting{mylibrary.py}

% integrate evaluate at (notice the left DOT!)
    \left .   ....   \right|_0^3
    $$ \int_0^3 x^2 dx = \left.\frac{x^3}{3} \right|_0^3 $$

% 4th root
    \sqrt[4]{...}

% font sizes
    \small
    \normalsize

% table example
    \setlength{\tabcolsep}{12pt}
    {\renewcommand{\arraystretch}{1.2}
    \begin{center}
    \begin{tabular}{l|c|c}
        ID & x & f(x)\\
        \hline
        A  & 5  & 25\\
        B  & 9  & 81\\
        C  & 1  & 1\\
        D  & 2  & 4
    \end{tabular}
    \end{center}

% approximately equals
    \approx

% verbatim
    \begin{verbatim}
    Some thing
         and leading space here
    \end{verbatim}

% much less than, much greater
    \ll \gg

% pound sterling
    \pounds

% actual backslashes
    \setminus
    \textbackslash

% set theory
    element \in
    subset \subseteq

% vertical spacing for align environment
% add in preamble.
    \addtolength{\jot}{0.5em}

% spacing
    \       % (space after backslash!) 	equivalent normal text space
    \quad 	% space equal to the current font size (= 18 mu)
    \, 	    % 3/18 of \quad (= 3 mu)
    \: 	    % 4/18 of \quad (= 4 mu)
    \; 	    % 5/18 of \quad (= 5 mu)
    \! 	    % -3/18 of \quad (= -3 mu)
    \qquad 	% twice \quad (= 36 mu)

% error for \url (missing package)
    \usepackage{url}
    \usepackage{hyperref}

% force figure placement
    \usepackage{float}
    \begin{figure}[H]

% increase item list separation
    \setlength\itemsep{1em}

% allow equations to span pagebreaks
    \allowdisplaybreaks

% evaluated at (vertical bar)
    \bigg|

% william stein errors
    \usepackage{amsmath}
    \DeclareMathOperator{\sss}{ss}

    \usepackage{amssymb}
    \newcommand{\R}{\mathbb{R}}

% cases example
    \[
    u(x) =
    \begin{cases}
        < 0 \implies \mbox{ maximum at} (x_0, f(x_o)) \\
        > 0 \implies \mbox{ minimum at} (x_0, f(x_o)) \\
        > 0 \implies \mbox{ minimum at} (x_0, f(x_o))
    \end{cases}
    \]

% colored box
\colorbox{gray}
{
    \parbox{\textwidth}
    {
        \textcolor{white}
        {
            Aside:\\
            \begin{align*}
                f(t) & = e^{g(t) } \\
                \frac{ df }{ dt } & = e^{g(t)}  \cdot \frac{ dg }{ dt }
            \end{align*}
        }
    }
}

% augmented matrix
    have defined: "augmat" snippet

% simultaneous equations
\begin{alignat*}{4}
    -2 x_1  & { }+{ } & x_2  & { }-{ } & x_3   & { }={ } & 4 \\
        x_1 & { }+{ } & 2x_2 & { }+{ } & 3x_3  & { }={ } & 13 \\
      3 x_1 & { }+{ } &      &           & x_3 & { }={ } & -1
\end{alignat*}

% force big brackets
     \big(
     \bigg(
     \Big(
     \Bigg(

% double bar x
    \overline{\overline{x}}

% therefore ("amssymb" package!)
    \therefore

% partial derivatives
    \frac{ \partial T }{ \partial n }

% plus minus
    \pm

% matrix styles
    matrix:  unbracketed matrix
    pmatrix: surrounded by (parentheses)
    bmatrix: surrounded by [square brackets]
    vmatrix: surrounded by |single vertical lines|
    Vmatrix: surrounded by ||double vertical lines||

% tag: use custom numbering of equations
\begin{align*}
    a & = b \\
    a & = b \tag{1}\\
    a & = b \\
\end{align*}

% example showing "nonumber" and "equation references" for the align environment
    \begin{align}
        a &=b \nonumber \\
        b &=c \label{eq:test}
    \end{align}
    Equation \eqref{eq:test} shows that $a=c$.

% numberthis custom command: for the align* environment
    \newcommand\numberthis{\addtocounter{equation}{1}\tag{\theequation}}
    \begin{align*}
        a & = someting \\
        a & = someting \numberthis \\
        a & = someting \\
    \end{align*}

% increase vertical line spacing for the array environment:
    \[ \arraycolsep=1.4pt\def\arraystretch{2.2}
    \begin{array}{lcl} \]

% maps to
    \to

% absolute bars (or determinant)
    \lvert M \rvert

% enumerate with letters
    \begin{enumerate}[(a)]

% enumerate with numbers
    \begin{enumerate}[1]

% enumerate with roman numerals
    \begin{enumerate}[(i)]

% use roman typeface in math mode
    $T_{\rm eff}$

% intersetion and union
    \cap and \cup

% ...
    \cdots

% make a vector tilde underneath a letter
    use ultisnip "vec"

% inline fractions
    \tfrac

% multiply
    \times

% Angstrom symbol
    ai texlive-science
    (siunix package)
    The lattice constant $a = \SI{2.0}{\angstrom}$ is ...
    if you only need the unit, use:
    Atomic distances are measured in \si{\angstrom}.

% shift a figure around
    \includegraphics[trim=2cm 0 0 0, scale=0.65]{ngts_planets.jpg}

% astro symbols
    % sun solar
    M_{\odot}

    % star
    M_{\star}

    % earth symbol/mass
    M_{\bigoplus}

    % moon
    \usepackage{wasysym}
    \leftmoon

% degree symbol
    $^{\circ}$

% dont have too much space around figures etc?
    \usepackage[nodisplayskipstretch]{setspace}

    dont know if this is any good too:
        % tweak the blank space around floating figures and tables
        % http://tex.stackexchange.com/questions/32614/vertical-spacing-in-floats
        % \setlength{\intextsep}{10mm}
        % \setlength{\belowcaptionskip}{-\baselineskip} \addtolength{\belowcaptionskip}{1.5ex}

% fancyhdr
    \today          % – the current day. ie: August 6, 2010
    \thepage        % – the current page number
    \thesection     % – the current section number
    \leftmark       % – current chapter name in uppercase. ie: CHAPTER 1. TITLE OF THE FIRST CHAPTER
    \rightmark      % – current section name in uppercase. ie: 1.1. TITLE OF THE FIRST SECTION

% typewriter font.
    \texttt

% compile latex (+ link in bibliography)
    % IN PREAMBLE: natbib package
        % \usepackage[square, numbers, comma, sort&compress]{natbib} <--- % numbered brackets [1], [2] etc
        \usepackage[authoryear, round, comma]{natbib}

    % IN DOCUMENT
        \section*{Bibliography}
        \bibliographystyle{unsrtnat}
        % \bibliographystyle{apalike}
        % \bibliographystyle{aa}
        \bibliography{references}

    % IN DOCUMENT
        and now we refer to that famous book \cite{griffeath2003}.

    % IN "references.bib"
        @book{griffeath2003,
        title={New Constructions in Cellular Automata},
        author={Griffeath, D. and Moore, C.},
        ...
        }

    % IN A SCRIPT:
        pdflatex notes.tex
        pdflatex notes.tex
        bibtex notes
        bibtex notes
        pdflatex notes.tex

% insert citation latex (authorea)
    \cite{myref}
    where {myref} is the first line of the entry in biblio.bib file

% cite author like "Hurley et al. ..."
    \citet{someref}

% cite page numbers
    \citep[pp. 10-12]{haswell_transiting_2010}

% color
    \usepackage{color}
    {\color{red} some text}

% multiply dot
    \cdot

% double subscript
    $a_{ij}$

% math mode!!!!
    use "$ ... $" or "\(" and "\)"

% best package
    "mathtools" is an extension of "amsmath"
    ai texlive-latex3
    \usepackage{mathtools}

% quit exit pdflatex compilation
    enter "x" at "?"

% tilde
    $\sim$

% put box around equation
    \boxed{}

% document order
    put packages before "begin document"!!!

% "\implies" not working
    usepackage{amsmath}
